% ============================================================
%  Manual do Usuário — Remote Firmware Flasher
%  Compilar com: tectonic manual.tex
% ============================================================
\documentclass[12pt,a4paper]{article}

% --- Codificação e língua ----------------------------------------
\usepackage[utf8]{inputenc}
\usepackage[T1]{fontenc}
\usepackage[brazil]{babel}

% --- Tipografia --------------------------------------------------
\usepackage{lmodern}
\usepackage{microtype}

% --- Layout ------------------------------------------------------
\usepackage[
  top=2.5cm, bottom=2.5cm,
  left=2.8cm, right=2.8cm,
  headheight=15pt
]{geometry}
\usepackage{fancyhdr}
\usepackage{titlesec}
\usepackage{parskip}
\setlength{\parindent}{0pt}
\setlength{\parskip}{6pt}
\usepackage{enumitem}

% --- Cor e links -------------------------------------------------
\usepackage{xcolor}
\definecolor{azulescuro}{RGB}{15,52,120}
\definecolor{azulmedio}{RGB}{40,100,180}
\definecolor{azulclaro}{RGB}{220,235,250}
\definecolor{cinzatexto}{RGB}{80,80,80}
\definecolor{amarelodest}{RGB}{255,243,205}
\definecolor{amarelob}{RGB}{180,130,0}

\usepackage[
  colorlinks=true,
  linkcolor=azulmedio,
  urlcolor=azulmedio,
  citecolor=azulmedio
]{hyperref}

% --- Tabelas -----------------------------------------------------
\usepackage{booktabs}
\usepackage{array}
\usepackage{colortbl}
\usepackage{tabularx}

% --- Caixas e destaques ------------------------------------------
\usepackage{tcolorbox}
\tcbuselibrary{skins,breakable}

\newtcolorbox{notabox}[1][Nota]{
  colback=amarelodest, colframe=amarelob,
  fonttitle=\bfseries\small, title=#1,
  left=4pt, right=4pt, top=3pt, bottom=3pt,
  breakable
}

\newtcolorbox{codigobox}{
  colback=gray!10, colframe=gray!40,
  fontupper=\ttfamily\small,
  left=8pt, right=4pt, top=4pt, bottom=4pt,
  breakable
}

% --- Código ------------------------------------------------------
\usepackage{listings}
\lstset{
  basicstyle=\ttfamily\footnotesize,
  keywordstyle=\color{azulmedio}\bfseries,
  backgroundcolor=\color{gray!8},
  frame=single,
  framerule=0.4pt,
  rulecolor=\color{gray!40},
  breaklines=true,
  showstringspaces=false
}

% --- Cabeçalho e rodapé ------------------------------------------
\pagestyle{fancy}
\fancyhf{}
\fancyhead[L]{\small\color{cinzatexto}Remote Firmware Flasher}
\fancyhead[R]{\small\color{cinzatexto}Manual do Usuário}
\fancyfoot[C]{\small\color{cinzatexto}\thepage}
\renewcommand{\headrulewidth}{0.4pt}
\renewcommand{\headrule}{\hbox to\headwidth{\color{azulmedio}\leaders\hrule height \headrulewidth\hfill}}

% --- Títulos de seção --------------------------------------------
\titleformat{\section}
  {\Large\bfseries\color{azulescuro}}
  {\thesection.}{0.6em}{}
  [\color{azulmedio}\rule{\linewidth}{1.2pt}]

\titleformat{\subsection}
  {\large\bfseries\color{azulmedio}}
  {\thesubsection}{0.6em}{}

% ================================================================
\begin{document}

% --- Capa --------------------------------------------------------
\begin{titlepage}
  \pagecolor{azulescuro}
  \color{white}
  \vspace*{2cm}

  \begin{center}
    {\fontsize{11}{13}\selectfont\scshape Residência Tecnológica --- UFPE/CIn}\\[0.8cm]
    \rule{\linewidth}{1.5pt}\\[0.6cm]
    {\fontsize{28}{34}\selectfont\bfseries Remote Firmware\\Flasher}\\[0.4cm]
    \rule{\linewidth}{1.5pt}\\[0.8cm]
    {\fontsize{15}{18}\selectfont Manual do Usuário}\\[0.3cm]
    {\fontsize{12}{15}\selectfont\color{azulclaro}%
      Sistemas Operacionais de Tempo Real}\\[2.5cm]
  \end{center}

  \begin{tcolorbox}[
    colback=white, colframe=white,
    coltext=azulescuro,
    width=0.75\textwidth,
    center, arc=3pt,
    left=10pt, right=10pt, top=6pt, bottom=6pt
  ]
    \small
    \begin{tabular}{@{}ll@{}}
      \textbf{Versão}          & 1.0 \\
      \textbf{Data}            & Março de 2026 \\
      \textbf{Repositório}     & \texttt{renatosfagundes/remote-flasher} \\
      \textbf{Plataforma}      & Windows 10/11 \\
    \end{tabular}
  \end{tcolorbox}

  \vfill
  {\small\itshape Gerado por \texttt{manual.tex} --- Compilado com Tectonic}
\end{titlepage}
\pagecolor{white}
\color{black}

% --- Sumário -----------------------------------------------------
\tableofcontents
\newpage

% === 1. Introdução ===============================================
\section{Introdução}

O \textbf{Remote Firmware Flasher} é uma aplicação desktop (PySide6) que permite programar placas Arduino remotamente, através de conexão VPN e SSH. Foi desenvolvido para a Residência Tecnológica do CIn/UFPE, permitindo que alunos gravem firmware nas placas do laboratório sem precisar estar fisicamente presente.

\subsection{Funcionalidades}

\begin{itemize}[nosep]
  \item Conexão VPN (SSTP) diretamente pela aplicação
  \item Upload e gravação de firmware (\texttt{.hex}) via SCP + \texttt{avrdude}
  \item Reset das placas via scripts PowerShell remotos
  \item Terminal serial para monitorar saída das placas em tempo real
  \item Terminal SSH para executar comandos no PC remoto
  \item Câmera ao vivo para visualizar as placas reagindo aos comandos
  \item Suporte a múltiplos PCs e placas do laboratório
  \item Configuração inicial automática (primeira execução)
  \item Criação automática da pasta remota do usuário
\end{itemize}

% === 2. Instalação ===============================================
\section{Instalação}

\subsection{Executável (recomendado)}

Basta baixar o arquivo \texttt{RemoteFlasher.exe} e executar. Não é necessário instalar Python ou qualquer dependência.

\subsection{Código-fonte}

Para executar a partir do código-fonte:

\begin{codigobox}
git clone https://github.com/renatosfagundes/remote-flasher.git\\
cd remote-flasher\\
pip install -r requirements.txt\\
python main.py
\end{codigobox}

\textbf{Dependências:} Python 3.9+, PySide6, paramiko, requests.

% === 3. Primeira Execução ========================================
\section{Primeira Execução}

Na primeira vez que você abrir o aplicativo, uma janela de configuração será exibida solicitando:

\begin{itemize}[nosep]
  \item \textbf{Seu nome} --- usado para criar sua pasta no PC remoto.
  \item \textbf{Pasta remota} --- preenchida automaticamente como \texttt{c:\textbackslash 2026\textbackslash <seu\_nome>}. Essa é a pasta no PC do laboratório onde seus arquivos \texttt{.hex} serão enviados.
\end{itemize}

Essas configurações são salvas em \texttt{\%APPDATA\%\textbackslash RemoteFlasher\textbackslash settings.json} e persistem entre sessões. Você pode limpar todas as configurações a qualquer momento pelo link ``Clear All Settings'' no rodapé da aba VPN.

A pasta remota é criada automaticamente no PC do laboratório na primeira operação que a utilize (flash, serial, reset, etc.).

% === 4. Conectando à VPN ========================================
\section{Conectando à VPN}

A aba \textbf{VPN} permite conectar-se à rede do laboratório:

\begin{itemize}[nosep]
  \item Nome da conexão: \texttt{VPN\_CIN} (já configurado)
  \item Endereço: \texttt{vpn.cin.ufpe.br}
  \item Protocolo: SSTP
  \item Usuário e senha: suas credenciais do CIn
\end{itemize}

\textbf{Passos:}
\begin{enumerate}[nosep]
  \item Insira seu usuário e senha na aba VPN
  \item (Opcional) Marque ``Remember me'' para salvar as credenciais localmente
  \item Clique em ``Connect VPN''
  \item Aguarde a conexão ser estabelecida (indicador verde)
\end{enumerate}

\begin{notabox}
Se a VPN ainda não estiver configurada no Windows, clique em ``Setup VPN'' para criar a conexão automaticamente.
\end{notabox}

\subsection{Limpar Configurações}

O link ``Clear All Settings'' no rodapé da aba VPN remove todas as configurações salvas (credenciais VPN, pasta remota) e fecha o aplicativo. Na próxima execução, o diálogo de primeira configuração será exibido novamente.

% === 5. Gravando Firmware ========================================
\section{Gravando Firmware (Flash)}

A aba \textbf{Flash} permite enviar e gravar firmware nas placas.

\textbf{Passos:}
\begin{enumerate}[nosep]
  \item Selecione o PC de destino (ex: PC 217)
  \item Selecione a placa (ex: Placa 01)
  \item Selecione a porta ECU (porta COM do Arduino no PC remoto)
  \item Clique em ``Browse'' e selecione seu arquivo \texttt{.hex}
  \item Clique em ``Flash''
\end{enumerate}

O aplicativo irá automaticamente:
\begin{enumerate}[nosep, label=\alph*)]
  \item Criar sua pasta remota (se não existir)
  \item Enviar o arquivo \texttt{.hex} via SCP
  \item Resetar a placa (se configurado)
  \item Gravar o firmware usando \texttt{avrdude}
  \item Reportar sucesso ou falha no log
\end{enumerate}

\begin{notabox}
Os arquivos \texttt{.hex} são gerados pelo Trampoline RTOS após a compilação. Exemplo: \texttt{lab02a\_fw.hex}
\end{notabox}

% === 6. Terminal Serial ==========================================
\section{Terminal Serial}

A aba \textbf{Serial} permite monitorar a saída serial das placas em tempo real através de SSH. O aplicativo executa o script \texttt{serialterm.py} no PC remoto e exibe a saída no terminal.

\textbf{Passos:}
\begin{enumerate}[nosep]
  \item Selecione o PC e a porta COM
  \item Configure o baudrate (padrão: 115200)
  \item Clique em ``Connect''
  \item A saída serial será exibida em tempo real
  \item Clique em ``Disconnect'' para encerrar
\end{enumerate}

% === 7. Terminal SSH =============================================
\section{Terminal SSH}

A aba \textbf{SSH Terminal} permite executar comandos arbitrários no PC remoto para depuração ou gerenciamento de arquivos.

\textbf{Passos:}
\begin{enumerate}[nosep]
  \item Selecione o PC de destino
  \item Opcionalmente, defina o diretório de trabalho remoto
  \item Digite o comando e pressione Enter ou clique em ``Run''
  \item A saída será exibida no terminal
\end{enumerate}

\begin{notabox}
O PC remoto usa Windows, então use comandos Windows (\texttt{dir}, \texttt{cd}, \texttt{copy}, etc.) e não comandos Linux (\texttt{ls}, \texttt{pwd}, \texttt{cp}, etc.).
\end{notabox}

% === 8. Reset de Placas ==========================================
\section{Reset de Placas}

A aba \textbf{Reset} permite enviar um sinal de reset para uma placa específica. O reset é feito executando scripts PowerShell (ex: \texttt{reset\_placa\_01.ps1}) no PC remoto.

\textbf{Passos:}
\begin{enumerate}[nosep]
  \item Selecione o PC e a placa
  \item Clique em ``Reset''
  \item Aguarde a confirmação no log
\end{enumerate}

% === 9. Câmera ===================================================
\section{Câmera ao Vivo}

Um painel de câmera pode ser exibido ao lado direito da aplicação, mostrando o feed de vídeo do laboratório em tempo real. Isso permite visualizar os LEDs das placas e confirmar visualmente que o firmware está funcionando.

\begin{itemize}[nosep]
  \item O painel aparece automaticamente nas abas Flash, Serial, SSH e Reset
  \item Na aba VPN, o painel de câmera fica oculto
  \item Você pode mostrar/ocultar o painel clicando no botão da câmera
\end{itemize}

% === 10. Configuração dos PCs ====================================
\section{Configuração dos PCs}

O arquivo \texttt{lab\_config.py} contém a configuração de todos os PCs e placas do laboratório. Os PCs disponíveis são:

\begin{center}
\begin{tabular}{llll}
  \toprule
  \textbf{PC} & \textbf{IP} & \textbf{Método Flash} & \textbf{Placas} \\
  \midrule
  PC 217 & 172.20.36.217 & avrdude  & 4 (com reset) \\
  PC 218 & 172.20.36.218 & avrdude  & 4 (sem reset) \\
  PC 220 & 172.20.36.220 & flash.py & 4 (com reset) \\
  PC 221 & 172.20.36.221 & avrdude  & 4 (sem reset) \\
  \bottomrule
\end{tabular}
\end{center}

Cada placa possui 4 portas ECU (Arduino), uma porta de reset e um script de reset associado. As portas COM específicas estão documentadas em \texttt{lab\_config.py}.

% === 11. Solução de Problemas ====================================
\section{Solução de Problemas}

\subsection{VPN não conecta}
\begin{itemize}[nosep]
  \item Verifique se suas credenciais estão corretas
  \item Tente o botão ``Setup VPN'' para recriar a conexão
  \item Verifique se outra VPN não está ativa
\end{itemize}

\subsection{Erro ao gravar firmware}
\begin{itemize}[nosep]
  \item Verifique se a porta COM está correta para a placa selecionada
  \item Tente resetar a placa antes de gravar novamente
  \item Verifique se o arquivo \texttt{.hex} é válido e não está corrompido
  \item Verifique se o baudrate do avrdude está correto (padrão: 57600)
\end{itemize}

\subsection{Terminal serial mostra caracteres estranhos}
\begin{itemize}[nosep]
  \item Verifique se o baudrate está correto (deve ser 115200 para Trampoline)
  \item Tente desconectar e reconectar
\end{itemize}

\subsection{Câmera não carrega}
\begin{itemize}[nosep]
  \item Verifique se o PC remoto tem o servidor de câmera ativo na porta 8080
  \item Tente acessar a URL da câmera diretamente no navegador
\end{itemize}

\subsection{Aplicativo fecha ao clicar em botões}
\begin{itemize}[nosep]
  \item Isso pode indicar um erro de configuração
  \item Tente executar a partir do código-fonte (\texttt{python main.py}) para ver mensagens de erro no terminal
\end{itemize}

\end{document}
